\documentclass[a4paper,11pt]{article}
        \usepackage[top=15mm,bottom=18mm,left=16mm,right=16mm]{geometry}
        \usepackage{fontspec}
        \usepackage{xeCJK}
        \setmainfont{Times New Roman}
        \setCJKmainfont{Yu Gothic}
        \setmonofont{Consolas}
        \setCJKmonofont{Yu Gothic}
        \usepackage{booktabs}
        \usepackage{longtable}
        \usepackage{tabularx}
        \usepackage{array}
        \usepackage{enumitem}
        \usepackage{hyperref}
        \usepackage{listings}
        \usepackage{xcolor}
        \usepackage{amsmath}
        \usepackage{amssymb}
        \hypersetup{
          colorlinks=true,
          linkcolor=blue,
          urlcolor=blue,
          pdftitle={ROS share / 相対パス解決},
          pdfauthor={Codex}
        }
        \lstset{
          basicstyle=\ttfamily\small,
          breaklines=true,
          columns=fullflexible,
          keepspaces=true,
          frame=single,
          framerule=0.2pt,
          rulecolor=\color{black},
          showstringspaces=false
        }
        \setlength{\parskip}{0.42em}
        \setlength{\parindent}{1em}
        \setlength{\emergencystretch}{3em}
        \renewcommand{\arraystretch}{1.15}
        \title{09. ROS share / 相対パス解決}
        \author{solar\_ws0129-main}
        \date{2026-07-29}
        \begin{document}
        \maketitle
        \tableofcontents
        \clearpage

        \section{対象}
        \begin{tabularx}{\linewidth}{>{\raggedright\arraybackslash}p{0.18\linewidth} X}
        \toprule
        項目 & 内容 \\
        \midrule
        ファイル & \path|mpc_solarcar/path_utils.py| \\
        ソースSHA-256 & \path|d2c07f2ccc796bf8d8d6149135b27123c8ecd873bc0b71c344028714a52c2a3b| \\
        種別 & Python \\
        区分 & config \\
        役割 & CWD、package share、repo root をまたいで path を実在ファイルへ解決する小さな基盤。 \\
        起動文脈 & Node 実行時の path ぶれを吸収する補助モジュール。 \\
        \bottomrule
        \end{tabularx}

        \section{このファイルを読む理由}
        CWD、package share、repo root をまたいで path を実在ファイルへ解決する小さな基盤。

        \begin{itemize}[leftmargin=1.6em]
        \item ament の package share があればそちらを優先する。
\item インストール後の launch/node 実行でも同じ relative path を使えるようにする。
        \end{itemize}

        \section{呼び出し関係}
        \subsection{どこから呼ばれるか}
        \begin{itemize}[leftmargin=1.6em]
        \item \path|mpc_node.py|
\item \path|gps_sim_node.py|
\item \path|solar_state_node.py|
        \end{itemize}

        \subsection{次に読むと理解が進むファイル}
        \begin{itemize}[leftmargin=1.6em]
        \item \path|mpc_solarcar/solar_profile.py|
        \end{itemize}

        \subsection{同一リポジトリ内の主要依存}
        \begin{itemize}[leftmargin=1.6em]
        \item 同一リポジトリ内の明示 import は少ない。
        \end{itemize}

        \section{ソース冒頭の把握}
        モジュール docstring または先頭意図の要約:

        モジュール docstring は明示されていない。

        \subsection{代表コード断片}
        \begin{lstlisting}[language=Python]
REPO_ROOT = Path(__file__).resolve().parents[1]


def resolve_path(path: str, default_subdir: str = '') -> str:
    """Resolve a path relative to CWD or package share.

    - If absolute, return as-is.
    - If exists relative to CWD, return it.
    - Otherwise, resolve under <pkg_share>/<default_subdir> (or directly under share).
    """
    if path is None:
        return path
    path = os.path.expanduser(str(path))
    if os.path.isabs(path):
        return path
    if os.path.exists(path):
        return path
    if get_package_share_directory is not None:
        pkg_share = get_package_share_directory(PKG_NAME)
    else:
        pkg_share = os.fspath(REPO_ROOT)
    if default_subdir:
\end{lstlisting}


        \section{主要構造の読み取り}
        主要関数は resolve\_path。

        \subsection{主要クラス・関数}
\begin{longtable}{p{0.07\linewidth} p{0.16\linewidth} p{0.25\linewidth} p{0.40\linewidth}}
\toprule
行 & 種別 & 名前 & 補足 \\
\midrule
14 & function & \path|resolve_path| & args: path, default\_subdir \\
\bottomrule
\end{longtable}

        \subsection{ファイルを上から読んだときの定義順}
            \begin{longtable}{p{0.07\linewidth} p{0.84\linewidth}}
            \toprule
            行 & 内容 \\
            \midrule
            4 & 例外処理を伴う try ブロックを実行する。 \\
10 & PKG\_NAME に 'mpc\_solarcar' の結果を代入する。 \\
11 & REPO\_ROOT に Path(\_\_file\_\_).resolve().parents[1] の結果を代入する。 \\
14 & 関数 resolve\_path を定義する。 \\
            \bottomrule
            \end{longtable}

        \subsection{import 群の詳細}
            \begin{longtable}{p{0.07\linewidth} p{0.33\linewidth} p{0.50\linewidth}}
            \toprule
            行 & import 文 & このファイルで読む理由 \\
            \midrule
            1 & \path|import os| & パス、環境変数、プロセス外部状態を扱うため。 このファイル内での主な使用位置は L23, L24, L26, L31, L34, L35, L36, L37。 \\
2 & \path|from pathlib import Path| & ファイルやディレクトリを安全に扱うため。 このファイル内での主な使用位置は L11。 \\
5 & \path|from ament_index_python.packages import get_package_share_directory| & ament\_index\_python.packages から get\_package\_share\_directory を読み込み、このファイルの処理を組み立てるため。 このファイル内での主な使用位置は L7, L28, L29。 \\
            \bottomrule
            \end{longtable}

        \section{このファイルを読む前に必要な基礎知識}
        次の章は、構文やROS用語を既知と仮定しないための説明である。

        \subsection{プログラム、プロセス、メモリ、オブジェクトを区別する}
ソースファイルはディスク上の文字列であり、それ自体は走っていない。Python実行可能プログラムがソースを読み、OSがその実行を一つのプロセスとして管理する。プロセスには仮想メモリ、開いているファイル、スレッド、終了コードなどが対応する。


\begin{lstlisting}
ソースファイル -> Pythonインタプリタが読み込む -> OS上のプロセス -> Pythonオブジェクトをメモリに生成 -> 関数やcallbackを実行
\end{lstlisting}


「メモリ上に生成する」とは、実行中プロセスが使う記憶領域に、その値の型、属性、参照関係を表す実体を用意することである。変数は実体そのものというより、そのオブジェクトを指す名前として理解するとPythonの代入が読みやすい。


\begin{lstlisting}[language=python]
node = MPCNode()
alias = node
# nodeとaliasは同じオブジェクトを参照する。
\end{lstlisting}


一つのプロセスには複数スレッドを持たせられる。スレッドは同じプロセスのメモリを共有するため、受信callbackと最適化callbackが同じself.zを書き換える場合は実行順と排他を検討する必要がある。

\paragraph{根拠資料}
\begin{itemize}[leftmargin=1.6em]
\item \href{https://docs.python.org/3/tutorial/classes.html}{Python公式チュートリアル: Classes}
\end{itemize}

\subsection{名前、代入、型、参照}
Pythonの代入文は、右辺を先に評価し、その結果のオブジェクトへ左辺の名前を結び付ける。`x = f()`では、まずfを呼び、その戻り値が得られてからxが更新される。


`float(...)`、`int(...)`、`bool(...)`は型名を呼び出して値を変換する。外部CSV、YAML、ROS parameterから来た値は型が期待どおりとは限らないため、このリポジトリでは明示変換が多い。


`None`は値が無いことを示す単一のオブジェクト、`math.nan`は浮動小数点値ではあるが有効な数値ではないことを示す。両者は用途が異なるため、`is None`と`math.isfinite`を使い分ける。

\paragraph{根拠資料}
\begin{itemize}[leftmargin=1.6em]
\item \href{https://docs.python.org/3/tutorial/classes.html}{Python公式チュートリアル: Classes}
\end{itemize}

\subsection{関数、引数、戻り値、スコープ}
`def`は関数本体をその場で実行する文ではない。関数オブジェクトを作り、その名前を現在の名前空間へ登録する。`f`は関数そのもの、`f()`はその関数を呼んだ結果である。


仮引数は関数定義側の受け取り名、実引数は呼出側から渡す値である。位置引数、キーワード引数、既定値付き引数、`*args`、`**kwargs`は渡し方が異なる。


\begin{lstlisting}[language=python]
def clip_speed(value, minimum=0.0, maximum=110.0):
    return min(maximum, max(minimum, value))

v = clip_speed(120.0, maximum=90.0)
\end{lstlisting}


関数内で作った通常の名前はローカルスコープに属する。メソッドが`self.z`を更新すればオブジェクトの状態が残るが、単に`z`を更新しただけならその呼出しのローカル値である。

\paragraph{根拠資料}
\begin{itemize}[leftmargin=1.6em]
\item \href{https://docs.python.org/3/tutorial/classes.html}{Python公式チュートリアル: Classes}
\end{itemize}

\subsection{module、package、import、console\_scripts}
Pythonファイルはmoduleとして読み込める。複数moduleをディレクトリにまとめたものがpackageである。`from .model import SolarCarModel`の先頭の点は、現在と同じpackage内のmodel moduleを指す相対importである。


import時にはファイルのトップレベル文が上から一度実行される。`def`や`class`は関数・クラスを登録するが、その本体の通常処理は呼び出すまで走らない。


setuptoolsの`console\_scripts`は、端末で使う実行可能名と`package.module:function`を対応付けるインストール時メタデータである。生成される実行用ラッパーはmoduleをimportし、指定関数を呼び、戻り値を終了コードとして扱う小さな入口である。


\begin{lstlisting}
ROS launchのexecutable名 -> install済みconsole script -> mpc_solarcar.mpc_nodeをimport -> main()を呼ぶ
\end{lstlisting}

\paragraph{根拠資料}
\begin{itemize}[leftmargin=1.6em]
\item \href{https://setuptools.pypa.io/en/latest/userguide/entry_point.html}{setuptools公式: Entry Points / Console Scripts}
\item \href{https://docs.python.org/3/tutorial/classes.html}{Python公式チュートリアル: Classes}
\end{itemize}

\subsection{例外、try/finally、with、資源解放}
例外は通常の戻り値とは別経路で異常を呼出元へ伝える。`try/except`は想定した異常を処理し、`finally`は成功・失敗にかかわらず後始末を行う。


`with open(...) as f:`はcontext managerを使い、ブロックを出るとファイルを閉じる。CSVやログの破損を避けるため、開いた資源の所有者と閉じる場所を明確にする。


`except Exception: pass`はノードを止めない利点がある一方、入力異常を隠して原因追跡を難しくする。安全に関係する値では、少なくとも頻度制限付き警告、異常カウンタ、fallback状態のpublishを検討する。



\subsection{freshness、filter、guard、fallback、fail-safe}
分散システムでは最後に受け取った値が現在も有効とは限らない。受信時刻とtimeoutからfreshnessを判定し、stale値を計画状態へ無条件に同期しない。


filterはnoiseと一時的な飛び値を抑えるが、遅れを生む。slew limitは指令変化率を制限する。安全guardはsolverのcost罰則とは別に、現在出力へ強制制約を適用する最後の防波堤である。


fallbackは失敗時の代替動作を事前に決める設計である。前回計画保持、物理に基づく決定論的入力、停止、低速制限などから、故障modeごとに選ぶ。fallback発生はstatusとlogへ残し、正常解と区別する。



\subsection{制御、状態、入力、モデル予測制御MPC}
制御対象の内部を表す状態をx、操作入力をu、外乱・予報をwとすると、離散モデルは`x[k+1] = f(x[k], u[k], w[k])`と書ける。ソーラーカーではSoC、電池温度、距離、速度などが状態候補、速度目標や駆動トルクが入力候補になる。


\[
\min_{u_0,\ldots,u_{N-1}} \sum_{k=0}^{N-1}\ell(x_k,u_k,w_k)+V_f(x_N)
\]

MPCは現在状態からNステップ先まで予測し、目的関数と制約を満たす入力系列を求める。ただし実際に適用するのは通常先頭入力だけで、次回は新しい実測状態から再び解く。これがreceding horizonである。


予測モデル、目的関数、制約、ホライズン、solver、初期値のどれかが変わると答えも変わる。「MPCを使う」だけでは仕様は決まらず、これらを単位付きで追う必要がある。

\paragraph{根拠資料}
\begin{itemize}[leftmargin=1.6em]
\item \href{https://docs.scipy.org/doc/scipy/reference/generated/scipy.optimize.minimize.html}{SciPy公式: scipy.optimize.minimize}
\end{itemize}



        \section{関数・クラスを上から順に解説}
        \subsection{L14 関数 resolve\_path}
            \begin{itemize}[leftmargin=1.6em]
              \item 行範囲: L14--L37
              \item 所属: module直下
              \item 定義: \path|resolve_path(path: str, default_subdir: str = '') -> str|
              \item docstring: Resolve a path relative to CWD or package share.

- If absolute, return as-is.
- If exists relative to CWD, return it.
- Otherwise, resolve under <pkg\_share>/<default\_subdir> (or directly under share).
              \item このブロックが直接呼ぶ主な関数・メソッド: exists, expanduser, fspath, get\_package\_share\_directory, isabs, join, startswith, str, strip
              \item 戻り値の要点: os.path.join(pkg\_share, path) / path / path / path
              \item 主なローカル代入名: path, pkg\_share, subdir
              \item 読み取る主なインスタンス属性: 特になし。
              \item 更新する主なインスタンス属性: 特になし。
              \item 明示的に送出する例外: 明示的raiseはない。
              \item 制御構造の規模: 条件分岐 6、ループ 0、try 0
            \end{itemize}
            \paragraph{この定義を読むためのPython構文}
            \begin{itemize}[leftmargin=1.6em]
            \item `->`以後は戻り値の型注釈であり、実行時の値を自動変換する命令ではない。
            \end{itemize}
            \paragraph{上から順の処理}
            \begin{enumerate}[leftmargin=1.8em]
            \item 条件 path is None を判定し、真なら内部処理を行う。
\item   path を返す。
\item path に os.path.expanduser(str(path)) の結果を代入する。
\item 条件 os.path.isabs(path) を判定し、真なら内部処理を行う。
\item   path を返す。
\item 条件 os.path.exists(path) を判定し、真なら内部処理を行う。
\item   path を返す。
\item 条件 get\_package\_share\_directory is not None を判定し、真なら内部処理を行う。
\item   pkg\_share に get\_package\_share\_directory(PKG\_NAME) の結果を代入する。
\item   上の条件が偽の場合:
\item   pkg\_share に os.fspath(REPO\_ROOT) の結果を代入する。
\item 条件 default\_subdir を判定し、真なら内部処理を行う。
\item   subdir に default\_subdir.strip('/\textbackslash{}\textbackslash{}') の結果を代入する。
\item   条件 path.startswith(subdir + os.sep) or path == subdir を判定し、真なら内部処理を行う。
\item     os.path.join(pkg\_share, path) を返す。
\item   os.path.join(pkg\_share, subdir, path) を返す。
\item os.path.join(pkg\_share, path) を返す。
            \end{enumerate}
            \paragraph{代表コード断片}
            \begin{lstlisting}[language=Python]
def resolve_path(path: str, default_subdir: str = '') -> str:
    """Resolve a path relative to CWD or package share.

    - If absolute, return as-is.
    - If exists relative to CWD, return it.
    - Otherwise, resolve under <pkg_share>/<default_subdir> (or directly under share).
    """
    if path is None:
        return path
    path = os.path.expanduser(str(path))
    if os.path.isabs(path):
        return path
    if os.path.exists(path):
        return path
    if get_package_share_directory is not None:
        pkg_share = get_package_share_directory(PKG_NAME)
    else:
        pkg_share = os.fspath(REPO_ROOT)
    if default_subdir:
        subdir = default_subdir.strip('/\\')
        if path.startswith(subdir + os.sep) or path == subdir:
            return os.path.join(pkg_share, path)
        return os.path.join(pkg_share, subdir, path)
    return os.path.join(pkg_share, path)
\end{lstlisting}











        \section{処理の流れ}
        \begin{enumerate}[leftmargin=1.7em]
        \item ソース中の主要関数を通じて処理を進める。
        \end{enumerate}

        \section{このファイルの位置づけ}
        この資料群では、\path|mpc_solarcar/path_utils.py| を
        \textbf{config}の中核として扱う。
        したがって、ここで扱う処理は単独で完結せず、
        前段の profile / launch / telemetry と後段の planner / logger / report へ繋がる。

        \end{document}
